% Generated by generate.py: do not edit.
\chapter{The Companion Volumes}
\section{Tome S1: Linear Algebra in Depth}
\subsection*{1.\ Spaces, Quotients and Duality}
\begin{results}
\item[1.1] Lemma (Steinitz exchange)
\item[1.2] Theorem (dimension)
\item[1.4] Theorem (quotients)
\item[1.7] Proposition (dual basis)
\item[1.8] Theorem (the double dual)
\item[1.10] Proposition (the trace pairing)
\item[1.12] Theorem (annihilators)
\item[1.13] Theorem (the transpose)
\item[1.14] Corollary (Fredholm alternative)
\item[1.16] Theorem (Oddtown; Berlekamp, 1969)
\item[1.17] Proposition (triangular criterion)
\item[1.19] Proposition (unisolvent points)
\end{results}
\subsection*{2.\ Determinants as Volume}
\begin{results}
\item[2.2] Theorem (exterior powers)
\item[2.5] Theorem (Cauchy--Binet)
\item[2.6] Corollary (Pythagoras for $k$-volumes)
\item[2.7] Theorem (Jacobi's complementary minors)
\item[2.8] Corollary (adjugate identity)
\item[2.9] Theorem (Desnanot--Jacobi identity)
\item[2.10] Corollary (Dodgson's condensation)
\item[2.11] Theorem (Lindstr\"om, 1973; Gessel--Viennot, 1985)
\item[2.14] Theorem (counting lozenge tilings)
\end{results}
\subsection*{3.\ Eigenvalues and Invariant Subspaces}
\begin{results}
\item[3.2] Theorem (triangularisation)
\item[3.3] Theorem (Cayley--Hamilton)
\item[3.4] Lemma (kernel lemma)
\item[3.5] Theorem (generalised eigenspaces)
\item[3.6] Corollary (diagonalisability)
\item[3.7] Theorem (Dunford--Jordan--Chevalley decomposition)
\item[3.9] Theorem (Jordan form of a nilpotent operator)
\item[3.10] Corollary (Jordan normal form)
\item[3.13] Proposition (the matrix exponential)
\end{results}
\subsection*{4.\ Inner Products and Projections}
\begin{results}
\item[4.2] Theorem (Cauchy--Schwarz)
\item[4.5] Theorem (Gram--Schmidt)
\item[4.6] Corollary ($QR$ factorisation)
\item[4.7] Theorem (orthogonal projection)
\item[4.8] Theorem (least squares)
\item[4.10] Theorem (Gauss--Markov)
\item[4.11] Theorem (Gram determinants)
\item[4.12] Theorem (Cauchy, 1841)
\item[4.15] Theorem (three-term recurrence)
\item[4.16] Theorem (zeros)
\item[4.17] Proposition (the Jacobi matrix)
\item[4.18] Theorem (Gauss quadrature)
\item[4.19] Theorem (Riesz representation)
\item[4.20] Proposition (the adjoint)
\item[4.21] Theorem (unitary operators)
\item[4.23] Corollary (Cholesky factorisation)
\end{results}
\subsection*{5.\ The Spectral Theorem and Its Relatives}
\begin{results}
\item[5.2] Theorem (spectral theorem)
\item[5.4] Theorem (Courant--Fischer)
\item[5.5] Theorem (Cauchy interlacing)
\item[5.6] Theorem (Weyl, 1912)
\item[5.7] Corollary (stability of eigenvalues)
\item[5.9] Theorem (Schur, 1923)
\item[5.10] Theorem (Horn, 1954)
\item[5.11] Theorem (Hoffman--Wielandt, 1953)
\end{results}
\subsection*{6.\ Singular Values}
\begin{results}
\item[6.1] Theorem (singular value decomposition)
\item[6.3] Theorem (polar decomposition)
\item[6.5] Theorem (Schmidt, Eckart--Young, Mirsky)
\item[6.8] Theorem (principal components)
\item[6.10] Theorem (orthogonal Procrustes)
\end{results}
\subsection*{7.\ Positive Matrices}
\begin{results}
\item[7.1] Theorem (Perron, 1907)
\item[7.3] Theorem (convergence of the power method)
\item[7.5] Theorem (Frobenius, 1912)
\item[7.8] Theorem (Wielandt)
\item[7.9] Corollary (ergodic theorem for Markov chains)
\end{results}
\subsection*{8.\ Graphs through Their Eigenvalues}
\begin{results}
\item[8.2] Theorem (Erd\H{o}s--R\'enyi--S\'os, 1966)
\item[8.4] Theorem (eigenvalues of strongly regular graphs)
\item[8.5] Theorem (Hoffman--Singleton, 1960)
\item[8.6] Theorem (Hoffman's ratio bound)
\item[8.8] Theorem (Graham--Pollak, 1971)
\item[8.9] Theorem (Huang, 2019)
\item[8.10] Proposition (Schwenk)
\end{results}
\subsection*{9.\ Symmetry Diagonalises}
\begin{results}
\item[9.1] Theorem (eigenvalues of circulants)
\item[9.7] Theorem (characters of a finite abelian group)
\item[9.8] Theorem (symmetry diagonalises)
\item[9.9] Theorem (Paley graphs)
\item[9.11] Theorem (distances in the cube)
\item[9.12] Theorem (Delsarte, 1973)
\item[9.13] Corollary (the linear programming bound)
\item[9.14] Corollary (Plotkin bound)
\item[9.16] Theorem (Dedekind)
\end{results}
\subsection*{10.\ Tensors and the Cost of Multiplying Matrices}
\begin{results}
\item[10.2] Proposition (rules of the Kronecker product)
\item[10.4] Proposition (the vec trick)
\item[10.5] Theorem (Sylvester's equation)
\item[10.7] Lemma (the slice test)
\item[10.12] Theorem (Strassen, 1969)
\item[10.13] Theorem (the exponent of matrix multiplication)
\item[10.14] Theorem (Schmidt decomposition)
\end{results}
\subsection*{11.\ Inequalities between Matrices}
\begin{results}
\item[11.2] Theorem (L\"owner--Heinz for the square root)
\item[11.3] Proposition (Schur complements)
\item[11.4] Proposition (Sylvester's determinant identity)
\item[11.5] Theorem (Hadamard's inequality)
\item[11.6] Corollary (Hadamard's volume inequality)
\item[11.7] Theorem (Fischer's inequality)
\item[11.8] Theorem (Minkowski's determinant inequality)
\item[11.9] Theorem (concavity of $\log\det$)
\item[11.11] Theorem (Schur, 1911)
\item[11.13] Theorem (Gershgorin, 1931)
\item[11.15] Theorem (Hilbert's inequality)
\end{results}
\subsection*{12.\ Linear Dynamics}
\begin{results}
\item[12.1] Lemma (a norm adapted to $A$)
\item[12.2] Theorem (convergence of powers)
\item[12.3] Theorem (Gelfand's formula, 1941)
\item[12.4] Theorem (stability)
\item[12.5] Theorem (Lyapunov, 1892)
\item[12.6] Theorem (Chebyshev acceleration)
\item[12.7] Theorem (heavy ball on quadratics)
\item[12.8] Theorem (Kalman's rank condition)
\item[12.10] Theorem (Maxwell, 1868)
\end{results}
\section{Tome S2: Geometry}
\subsection*{1.\ Euclid Revisited}
\begin{results}
\item[1.1] Proposition (three tests)
\item[1.2] Lemma (chords of the unit circle)
\item[1.3] Proposition (the orthocentre)
\item[1.4] Theorem (Euler, 1765)
\item[1.5] Theorem (the nine-point circle)
\item[1.6] Proposition (the incentre)
\item[1.7] Theorem (Ceva)
\item[1.10] Theorem (Napoleon's theorem)
\item[1.11] Theorem (Morley, 1899)
\end{results}
\subsection*{2.\ Transformations}
\begin{results}
\item[2.1] Theorem (isometries are affine)
\item[2.2] Theorem (three reflections suffice)
\item[2.3] Theorem (classification of plane isometries)
\item[2.5] Theorem (Leonardo's theorem)
\item[2.7] Theorem (the seven friezes)
\item[2.8] Theorem (the crystallographic restriction)
\item[2.9] Proposition (the five lattice types)
\item[2.10] Theorem (Fagnano)
\end{results}
\subsection*{3.\ Inversion}
\begin{results}
\item[3.2] Proposition (distances)
\item[3.3] Theorem (inversion and circles)
\item[3.5] Theorem (Ptolemy's inequality)
\item[3.6] Proposition (inversive distance is invariant)
\item[3.7] Lemma (concentric position)
\item[3.8] Theorem (Steiner's porism)
\item[3.9] Theorem (Pappus)
\item[3.10] Theorem (Descartes 1643, Soddy 1936; Lagarias, Mallows and Wilks 2002)
\item[3.11] Proposition (the swap rule)
\item[3.12] Corollary (integral gaskets)
\item[3.13] Theorem (M\"obius maps)
\item[3.14] Theorem (Feuerbach, 1822)
\end{results}
\subsection*{4.\ Projective geometry}
\begin{results}
\item[4.3] Theorem (Duality)
\item[4.5] Proposition (Frames)
\item[4.6] Theorem (Desargues)
\item[4.7] Theorem (Pappus)
\item[4.12] Theorem (Harmonic construction)
\item[4.13] Theorem (Fundamental theorem of projective geometry)
\end{results}
\subsection*{5.\ Conics}
\begin{results}
\item[5.4] Theorem (Normal form)
\item[5.5] Proposition (Focus and directrix)
\item[5.6] Theorem (Chasles and Steiner)
\item[5.8] Theorem (Pascal)
\item[5.11] Theorem (Poncelet's porism for two circles)
\end{results}
\subsection*{6.\ Quadratic forms and quadrics}
\begin{results}
\item[6.1] Proposition (Orthogonal bases)
\item[6.2] Theorem (Sylvester's law of inertia)
\item[6.5] Theorem (The Segre quadric)
\item[6.8] Theorem (Witt cancellation)
\item[6.9] Corollary (Witt decomposition)
\end{results}
\subsection*{7.\ Plane curves}
\begin{results}
\item[7.4] Theorem (B\'ezout)
\item[7.9] Theorem (Cayley--Bacharach for cubics)
\item[7.10] Corollary (Pascal, again)
\end{results}
\subsection*{8.\ Cubic curves}
\begin{results}
\item[8.4] Theorem (associativity)
\item[8.5] Proposition (Weierstrass form)
\end{results}
\subsection*{9.\ Lattices and sphere packings}
\begin{results}
\item[9.3] Lemma (Blichfeldt)
\item[9.4] Theorem (Minkowski)
\item[9.5] Theorem (Fermat)
\end{results}
\subsection*{10.\ Non-Euclidean geometry}
\begin{results}
\item[10.1] Theorem (Girard, 1629)
\item[10.2] Proposition (spherical law of cosines)
\item[10.4] Proposition (isometries)
\item[10.6] Theorem (hyperbolic distance)
\item[10.7] Proposition (many parallels)
\item[10.8] Theorem (Gauss--Bonnet for triangles)
\item[10.10] Proposition (the disk model)
\end{results}
\subsection*{11.\ Convex bodies}
\begin{results}
\item[11.1] Lemma (Radon, 1921)
\item[11.2] Theorem (Helly)
\item[11.4] Theorem (Buffon, Barbier)
\item[11.5] Theorem (Cauchy's formula)
\item[11.6] Corollary (Barbier, 1860)
\item[11.7] Theorem (Brunn--Minkowski)
\item[11.8] Corollary (isoperimetric inequality)
\end{results}
\subsection*{12.\ Curvature}
\begin{results}
\item[12.1] Theorem (curvature determines the curve)
\item[12.2] Theorem (four vertex theorem)
\item[12.5] Theorem (Gauss--Bonnet for convex surfaces)
\item[12.6] Theorem (Euler, 1758)
\item[12.7] Theorem (Descartes)
\end{results}
\section{Tome S3: A Book of Conjectures}
\subsection*{1.\ How Conjectures Are Born and Die}
\begin{results}
\item[1.1] Theorem (Landau, 1905)
\item[1.3] Theorem (P\'olya, 1919)
\end{results}
\subsection*{2.\ Twin Primes and Prime Gaps}
\begin{results}
\item[2.4] Theorem (Erd\H{o}s, 1940)
\item[2.5] Theorem (Westzynthius, 1931)
\end{results}
\subsection*{3.\ Goldbach}
\begin{results}
\item[3.1] Theorem (Schnirelmann's inequality)
\item[3.6] Theorem (Schnirelmann, 1930)
\end{results}
\subsection*{4.\ The Riemann Hypothesis}
\begin{results}
\item[4.1] Theorem (Gronwall, 1913)
\item[4.2] Proposition (Alaoglu--Erd\H{o}s)
\item[4.4] Proposition (Lagarias, the easy direction)
\end{results}
\subsection*{5.\ Zeta and Random Matrices}
\begin{results}
\item[5.3] Theorem (Wigner, 1955)
\item[5.4] Proposition (the Wigner surmises)
\item[5.7] Theorem (mean square to the right of the line)
\end{results}
\subsection*{6.\ Collatz}
\begin{results}
\item[6.4] Theorem (Terras, 1976; Everett, 1977)
\item[6.5] Corollary (the exact random model)
\item[6.6] Lemma (Chernoff bound)
\item[6.7] Theorem (Terras, 1976)
\item[6.8] Proposition (the cycle equation)
\item[6.9] Lemma (Farey neighbours)
\item[6.10] Theorem (cycles are long)
\end{results}
\subsection*{7.\ Hadamard Matrices}
\begin{results}
\item[7.1] Theorem (Hadamard, 1893)
\item[7.3] Proposition (Sylvester, 1867)
\item[7.6] Theorem (Paley I, 1933)
\item[7.7] Theorem (Paley II, 1933)
\item[7.8] Theorem (Williamson, 1944)
\end{results}
\subsection*{8.\ Points, Distances and Families of Sets}
\begin{results}
\item[8.1] Theorem (Moser and Moser, 1961; Isbell, 1950)
\item[8.2] Theorem (Spencer, Szemer\'edi and Trotter, 1984; Sz\'ekely's proof)
\item[8.4] Theorem (Erd\H os and Szekeres, 1935)
\item[8.8] Theorem (Gilmer's method)
\end{results}
\subsection*{9.\ The abc Conjecture}
\begin{results}
\item[9.2] Theorem (Stothers, 1981; Mason, 1984)
\item[9.3] Corollary (Fermat's last theorem for polynomials)
\item[9.6] Theorem (abc implies the Fermat--Catalan conjecture)
\item[9.9] Theorem (Mordell)
\end{results}
\subsection*{10.\ Birch and Swinnerton-Dyer}
No named results; 8 exercises.
\subsection*{11.\ Randomness in the Integers}
No named results; 8 exercises.
\subsection*{12.\ Open Doors in Algebra}
\begin{results}
\item[12.6] Theorem (after Breusch and Smyth)
\item[12.8] Theorem (Wang, 1980)
\end{results}
